% ============================================================================
% Manuscript fixes for "Revising pyBuildingEnergy for an Australian Apartment"
% Each block names the location, what to change, and LaTeX to paste.
% All numbers below were recomputed from the saved hourly outputs; none of the
% existing results in the paper change.
% Group A = required for consistency; Group B = strengthens the paper;
% Group C = small corrections.
% ============================================================================


% ----------------------------------------------------------------------------
% A1. Section 4.4 -- DELETE the duplicated paragraph
% ----------------------------------------------------------------------------
% Delete the whole paragraph that begins
%   "The combined demand change was nearly additive in pyBuildingEnergy. The
%    interaction, defined as ... approximately -0.1 kWh. The corresponding
%    interaction calculated from the displayed EnergyPlus totals ..."
% The paragraph before it already reports both interactions from unrounded
% outputs; this one repeats it and contradicts it ("displayed totals").


% ----------------------------------------------------------------------------
% A2. Abstract -- state that the headline figures use the diagnostic construction
% ----------------------------------------------------------------------------
% REPLACE: "Combined implementation and boundary-condition changes reduce
%           annual sensible demand from 2408.74 to 131.76 kWh"
% WITH:
For the diagnostic zero-capacity construction, combined implementation and
boundary-condition changes reduce annual sensible demand from 2408.74 to
131.76~kWh
% Make the same qualification in the first sentence of the second paragraph
% of the Conclusion ("Stage 1 produced the largest improvement. For the
% diagnostic construction, annual sensible demand decreased ...").


% ----------------------------------------------------------------------------
% A3. Section 3.1 (or a row in Table 1) -- operating schedules
% ----------------------------------------------------------------------------
% ADD after the sentence introducing Table 1:
Conditioning is intermittent. The heating setpoint of 18\,\textdegree C applies
from 19:00 to 23:00 on weekdays and from 18:00 to 23:00 at weekends, with a
15\,\textdegree C setback at other times. The cooling setpoint of
26\,\textdegree C applies from 17:00 to 23:00 on weekdays and from 09:00 to
23:00 at weekends, with a 28\,\textdegree C setback otherwise. Designed
ventilation operates continuously. Both engines use the same weekday and
weekend schedules.


% ----------------------------------------------------------------------------
% A4. Section 3.5.1 -- say which factor the fixed case uses
% ----------------------------------------------------------------------------
% ADD at the end of the first paragraph (after "...regardless of the actual
% direction of sunlight."):
The fixed-transmittance configurations in this study use a multiplier of 1;
the constant 0.9 option is not used.


% ----------------------------------------------------------------------------
% A5. Section 3.5 (end of the introductory paragraph) -- code defaults
% ----------------------------------------------------------------------------
% ADD:
In the released code, both modifications are enabled by default; the
fixed-method configurations disable them explicitly.


% ----------------------------------------------------------------------------
% A6. Appendix D3 -- external shading
% ----------------------------------------------------------------------------
% ADD at the end of the first paragraph of D3:
No external shading is represented in either engine. The case documentation
describes a 0.25\,m overhang above the windows; it is excluded from both input
models so that the reference contains no unmatched shading surface.


% ----------------------------------------------------------------------------
% A7. Appendix D4 -- source of the gain levels
% ----------------------------------------------------------------------------
% REPLACE: "Corrected internal-gain design levels are 84 W for occupants, 60 W
%           for appliances and zero for lighting, on the shared schedules."
% WITH:
Internal-gain design levels are the EN~16798-1 residential-apartment values of
4.2, 3.0 and 0~W/m\textsuperscript{2} for occupants, appliances and lighting,
giving 84, 60 and 0~W over the 20~m\textsuperscript{2} floor area, on the
shared hourly schedules.


% ----------------------------------------------------------------------------
% A8. Appendix A1 -- multiplier value
% ----------------------------------------------------------------------------
% REPLACE: "the multiplier is approximately 7.33"
% WITH:
the multiplier is 7.335
% (The saved gain totals give exactly 5341.22 / 728.18 = 7.335.)


% ----------------------------------------------------------------------------
% A9. Section 4.5 -- historical reproduction and verification
% ----------------------------------------------------------------------------
% The claim that the historical case "reproduced the earlier results within
% their reported precision" could not be traced to a specific earlier record.
% EITHER cite the record it reproduces, OR replace the sentence with:
These values provide continuity with the configuration used in earlier
versions of the analysis; they are not used in the staged comparisons.

% REPLACE the test-count sentence with (counts after the input-handling fixes):
The test suite recorded 261 passed and 17 skipped checks, including 12
regression checks for input handling. The skipped checks compare against
historical development branches, are not executable from the archived
release, and are not counted as verification evidence. A fresh annual run of
all twelve main cases reproduced the archived hourly conditioning loads with a
maximum difference of 0~W.


% ----------------------------------------------------------------------------
% A10. Appendix A3 -- explicit coldest month (true after the code fix)
% ----------------------------------------------------------------------------
% No wording change is needed once the engine fix is applied; the explicit
% month is now read from building_parameters.climate_parameters.coldest_month.
% If you publish WITHOUT the code fix, replace "unless a valid explicit month
% is supplied" with "unless building_parameters.coldest_month is supplied".


% ============================================================================
% GROUP B -- strengthens the paper
% ============================================================================

% ----------------------------------------------------------------------------
% B1. Discussion -- new paragraph on the remaining heating discrepancy
%     (insert after the paragraph on Stage 1, before the Stage 2 paragraph)
% ----------------------------------------------------------------------------
The remaining heating underprediction can be located more precisely.
Table~\ref{tab:winter-balance} compares the energy-balance terms of the repaired
configuration with fixed methods over the 430 June--August hours in which either
engine supplied heating. Both engines maintained the same mean operative
temperature (17.51 and 17.49\,\textdegree C), and their exterior-wall and
outdoor-air terms agreed within about 1~kWh. The difference arose at the shared
boundaries: pyBuildingEnergy received 19.6~kWh more heat from the five
neighbouring spaces, which accounts for the 16.1~kWh winter heating shortfall,
partly offset by window exchange. The shortfall was distributed across all
heated hours rather than concentrated at the start of each heating period. The
comparison identifies where the engines diverge but not the mechanism; the
treatment of interior surface exchange, which differs between the engines
(Appendix~D3), is the most direct candidate for a controlled test. The
residual therefore lies in the same shared boundaries that dominate the
apartment's conductance and carry the assumed 20\,\textdegree C neighbouring
temperature, rather than in the glazing and exterior-convection mechanisms
evaluated in Stage~2.

\begin{table}[htbp]
\centering
\caption{Energy-balance terms during winter heated hours (June--August, 430~h),
repaired configuration with fixed methods. Energy in kWh; gains positive.}
\label{tab:winter-balance}
\begin{tabular}{lrrr}
\toprule
Term & pyBuildingEnergy & EnergyPlus & Difference \\
\midrule
Heating supplied                     &  78.1 &  94.2 & $+16.1$ \\
Heat gained from neighbouring spaces & 160.6 & 141.0 & $-19.6$ \\
Exterior wall                        & $-53.3$ & $-54.2$ & $-0.9$ \\
Outdoor air (loss)                   & 190.8 & 189.6 & $-1.3$ \\
Window, net                          & $-39.0$ & $-32.2$ & $+6.8$ \\
\bottomrule
\end{tabular}
\end{table}
% Difference = EnergyPlus minus pyBuildingEnergy.


% ----------------------------------------------------------------------------
% B2. Section 3.6.1 -- implied air-change rates and the ventilation share
% ----------------------------------------------------------------------------
% ADD after "...and 14 provides an additional sensitivity case.":
Applied to the 13.5~m\textsuperscript{2} outdoor-exposed envelope and 54~m\textsuperscript{3}
volume, these values correspond to $n_{50}$ = 1.00, 2.75 and 3.50~h$^{-1}$,
well below the whole-dwelling averages reported in the Australian surveys.
The continuous designed ventilation also dominates outdoor-air exchange: its
conductance of 48.4~W/K compares with mean infiltration conductances of 0.76,
2.09 and 2.66~W/K, so infiltration contributes 1.5--5.2\% of outdoor-air
conductance in the tested scenarios.

% And in Section 4.3, REPLACE "These effects are conditional on the
% apartment's small outdoor-exposed envelope and continuous designed
% ventilation of 0.040 m3/s." WITH:
These effects are conditional on the apartment's small outdoor-exposed
envelope and continuous designed ventilation of 0.040~m\textsuperscript{3}/s,
under which infiltration provides at most 5.2\% of outdoor-air conductance.


% ----------------------------------------------------------------------------
% B3. Section 4.3 -- airflow timing isolated (pyBuildingEnergy side)
% ----------------------------------------------------------------------------
% REPLACE the two sentences beginning "However, this experiment changes both
% mean airflow and its hourly distribution..." WITH:
An additional pyBuildingEnergy run used constant airflow scaled to the weather
multiplier's annual mean (0.853), so that mean airflow matched the
weather-dependent case. Weather-dependent airflow then increased heating by
2.79~kWh (120.73 against 117.94~kWh) and cooling by 0.32~kWh, isolating the
effect of timing in this engine. A matched EnergyPlus reference was not run.
% NOTE: if you add this sentence, add the run to Table 12 (Appendix C4) as
% "Mean-matched constant flow, q50 = 9.387 (11 x 0.853), D, C, Both, --" and
% keep the run files. For a matched EnergyPlus figure, one more reference is
% needed (infiltration design flow x 0.853, constant schedule).


% ----------------------------------------------------------------------------
% B4. Table 6 -- add hourly RMSE, and note zero cooling
% ----------------------------------------------------------------------------
% Add a final column "RMSE (W)" with values:
%   Stage 3 reference                15.02
%   Higher leakage only              15.76
%   Material-derived construction    19.39
%   Combined input scenario          20.29
% ADD to the paragraph after Table 6:
Hourly agreement also deteriorated with material-derived construction:
total-load RMSE increased from 15.02 to 19.39~W, and to 20.29~W for the
combined scenario. Neither engine predicted any cooling demand in these
configurations, so their total bias equals the heating bias.


% ----------------------------------------------------------------------------
% B5. Section 4.2 -- justification for retaining both Stage 2 methods
% ----------------------------------------------------------------------------
% REPLACE "Both methods were retained for Stage 3 as a consistent tested
% configuration, rather than as a configuration that minimised every error
% measure." WITH:
Both methods were retained for Stage~3 on physical grounds rather than
because they minimised every error measure: transmission through glazing
depends on incidence angle, and the EnergyPlus reference itself uses
wind-dependent exterior convection.


% ----------------------------------------------------------------------------
% B6. Glazing description -- choose ONE option (Section 3.5.1 or Appendix B1)
% ----------------------------------------------------------------------------
% A U-value of 5.4 W/(m2 K) normally indicates single glazing, whereas the
% angular correction uses the two-pane curve.
% OPTION 1 (keep runs; justify):
The two-pane coefficients are an assumed optical response; the window U-value
of 5.4~W/(m\textsuperscript{2}\,K) is more typical of single glazing, so the
angular correction should be regarded as an assumption rather than a property
of the installed windows.
% OPTION 2 (stronger): add a one-pane sensitivity run. The code already
% provides one-pane coefficients (glazing_panes = 1). One ISO run, no new
% EnergyPlus reference needed.


% ============================================================================
% GROUP C -- small corrections
% ============================================================================

% C1. Section 3.2: REPLACE "Wind direction, shelter from nearby buildings, and
%     facade pressure coefficients are kept as is."
%     WITH:
Wind direction, shelter from nearby buildings and facade pressure
coefficients are not modelled.

% C2. Section 3.3 vs Appendix A5 -- make the short-wave wording agree.
%     REPLACE in Section 3.3 "It also includes absorbed external short-wave
%     radiation." WITH:
It is net of absorbed external short-wave radiation, which is reported
separately.

% C3. Leakage values and the archived age lookup -- ADD to Appendix C1:
The archived lookup assigns $q_{50}$ = 11 to dwellings built from 2006 and
14 to those built from 1991 to 2005; the selected scenario values coincide
with these entries but are used here as explicit illustrative inputs.

% C4. References [15] and [16] are the same CSIRO article. Delete one and
%     update the citations.

% C5. Lower-case "Solar transmittance" mid-sentence (Sections 2.1, 4.2, 5, 6).

% C6. Section 3.1: "11.88 m2. of opaque wall and two windows with a combined
%     area of 1.62 m2 The remaining" -> move the full stop to after "1.62 m2".
